% math0 manuscript. Build with: latexmk -pdf paper.tex
\documentclass[11pt]{amsart}
\usepackage[margin=1.1in]{geometry}
\usepackage{amsmath,amssymb,amsthm,mathtools}
\usepackage{url}
\usepackage{hyperref}
\usepackage{xcolor}
\usepackage{listings}
\usepackage{newunicodechar}

\lstdefinelanguage{lean}{
  keywords={theorem,lemma,def,abbrev,example,fun,by,have,show,from,let,in,if,then,else,match,with,
    namespace,end,open,import,section,variable,instance,structure,where,calc,obtain,intro,exact,
    refine,rcases,rintro,simp,omega,decide,norm_num,linarith,nlinarith,positivity,induction,cases,
    private},
  sensitive=true, comment=[l]{--}, morecomment=[s]{/-}{-/}, string=[b]"}
\lstset{language=lean, basicstyle=\ttfamily\footnotesize, columns=fullflexible, keepspaces=true, breaklines=true,
  keywordstyle=\color{blue!60!black}\bfseries, commentstyle=\color{gray}, frame=single,
  extendedchars=true, inputencoding=utf8,
  literate={ℕ}{{$\mathbb{N}$}}1 {ℤ}{{$\mathbb{Z}$}}1 {∀}{{$\forall$}}1 {↔}{{$\leftrightarrow$}}1
    {¬}{{$\neg$}}1 {∧}{{$\wedge$}}1 {∨}{{$\vee$}}1 {→}{{$\to$}}1 {≤}{{$\le$}}1 {<}{{$<$}}1}
% Unicode commonly found in Lean sources.
\newunicodechar{ℕ}{\ensuremath{\mathbb{N}}}
\newunicodechar{ℤ}{\ensuremath{\mathbb{Z}}}
\newunicodechar{ℚ}{\ensuremath{\mathbb{Q}}}
\newunicodechar{ℝ}{\ensuremath{\mathbb{R}}}
\newunicodechar{ℂ}{\ensuremath{\mathbb{C}}}
\newunicodechar{∀}{\ensuremath{\forall}}
\newunicodechar{∃}{\ensuremath{\exists}}
\newunicodechar{→}{\ensuremath{\to}}
\newunicodechar{←}{\ensuremath{\leftarrow}}
\newunicodechar{↔}{\ensuremath{\leftrightarrow}}
\newunicodechar{↦}{\ensuremath{\mapsto}}
\newunicodechar{¬}{\ensuremath{\neg}}
\newunicodechar{∧}{\ensuremath{\wedge}}
\newunicodechar{∨}{\ensuremath{\vee}}
\newunicodechar{≤}{\ensuremath{\le}}
\newunicodechar{≥}{\ensuremath{\ge}}
\newunicodechar{≠}{\ensuremath{\ne}}
\newunicodechar{∈}{\ensuremath{\in}}
\newunicodechar{∉}{\ensuremath{\notin}}
\newunicodechar{⊆}{\ensuremath{\subseteq}}
\newunicodechar{∩}{\ensuremath{\cap}}
\newunicodechar{∪}{\ensuremath{\cup}}
\newunicodechar{∑}{\ensuremath{\sum}}
\newunicodechar{∏}{\ensuremath{\prod}}
\newunicodechar{⟨}{\ensuremath{\langle}}
\newunicodechar{⟩}{\ensuremath{\rangle}}
\newunicodechar{·}{\ensuremath{\cdot}}
\newunicodechar{∘}{\ensuremath{\circ}}
\newunicodechar{⁻}{\ensuremath{^{-}}}
\newunicodechar{¹}{\ensuremath{^{1}}}
\newunicodechar{²}{\ensuremath{^{2}}}
\newunicodechar{₀}{\ensuremath{_{0}}}
\newunicodechar{₁}{\ensuremath{_{1}}}
\newunicodechar{₂}{\ensuremath{_{2}}}
\newunicodechar{α}{\ensuremath{\alpha}}
\newunicodechar{β}{\ensuremath{\beta}}
\newunicodechar{ε}{\ensuremath{\varepsilon}}
\newunicodechar{δ}{\ensuremath{\delta}}
\newunicodechar{π}{\ensuremath{\pi}}
\newunicodechar{σ}{\ensuremath{\sigma}}
\newunicodechar{φ}{\ensuremath{\varphi}}
\newunicodechar{∞}{\ensuremath{\infty}}
\newunicodechar{⌊}{\ensuremath{\lfloor}}
\newunicodechar{⌋}{\ensuremath{\rfloor}}
\newunicodechar{⌈}{\ensuremath{\lceil}}
\newunicodechar{⌉}{\ensuremath{\rceil}}
\newunicodechar{∣}{\ensuremath{\mid}}
\newunicodechar{≡}{\ensuremath{\equiv}}
\newunicodechar{×}{\ensuremath{\times}}
\newunicodechar{⧸}{\ensuremath{/}}
\newunicodechar{‖}{\ensuremath{\|}}
\newunicodechar{√}{\ensuremath{\surd}}
\newunicodechar{∅}{\ensuremath{\emptyset}}
\newunicodechar{⦃}{\ensuremath{\{\!|}}
\newunicodechar{⦄}{\ensuremath{|\!\}}}
\newunicodechar{𝓝}{\ensuremath{\mathcal{N}}}
\newunicodechar{≫}{\ensuremath{\gg}}
\newunicodechar{≪}{\ensuremath{\ll}}

\theoremstyle{plain}
\newtheorem{theorem}{Theorem}
\newtheorem{lemma}[theorem]{Lemma}
\newtheorem{proposition}[theorem]{Proposition}
\newtheorem{corollary}[theorem]{Corollary}
\theoremstyle{definition}
\newtheorem{definition}[theorem]{Definition}
\newtheorem{remark}[theorem]{Remark}

\emergencystretch=3em
% Lean identifiers in running text (breakable at dots and underscores).
\expandafter\def\expandafter\UrlBreaks\expandafter{\UrlBreaks\do\_\do\.}
\DeclareUrlCommand\decl{\urlstyle{tt}}
\newcommand{\N}{\mathbb{N}}
\newcommand{\Z}{\mathbb{Z}}
\newcommand{\Lev}{\mathrm{Lev}}
\newcommand{\last}{\operatorname{last}}
\newcommand{\len}[1]{\lvert #1\rvert}
\newcommand{\A}{A}

\title[The Golomb triangle A113676 and Levine's sequence A014644]{The row lengths and row ends of the Golomb triangle OEIS A113676 are given by Levine's sequence A014644}
\author{math0}
\address{Produced by Claude (an AI model developed by Anthropic) in the math0 project}
\date{October 20, 2026}

\begin{document}
\begin{abstract}
Cut Golomb's sequence with one of its two $2$'s removed, $1,2,3,3,4,4,4,5,5,5,\dots$, into consecutive rows so that the first row is $(1)$ and each further row has as many terms as the sum of the previous row; the row lengths form OEIS A113676.
We prove both conjectures that Floor van Lamoen and Paul D.\ Hanna recorded in that entry in 2005: for $n\ge2$ the length of row $n$ equals $A(n+2)-A(n+1)$, and the last term of row $n$ equals $A(n+1)$, where $A$ is Levine's sequence A014644 (and row $1$ ends in $A(1)=1$).
The proof rests on the fact, also observed by Claes and Miyamoto, that from the second row on the rows of Levine's array are initial segments of Golomb's sequence; a telescoping identity and an induction then locate the rows of the triangle exactly.
To our knowledge the result is new, though elementary; the OEIS entry still lists both statements as conjectures.
The theorem is verified in Lean~4 against a frozen, fidelity-reviewed statement that math0 formalized from the OEIS entry, whose definitions are checked by the kernel against the listed OEIS terms.
\end{abstract}
\maketitle

\section{Introduction}

\emph{Golomb's sequence} $G$ (OEIS A001462 \cite{A001462}; Golomb \cite{Go}, Guy \cite[E25]{Guy}) is the unique nondecreasing sequence of positive integers with $G(1)=1$ in which every $n\ge1$ occurs exactly $G(n)$ times:
\[
  G = 1,2,2,3,3,4,4,4,5,5,5,6,6,6,6,7,\dots
\]
\emph{Levine's array} A014643 \cite{A014643} starts with the row $(1,1)$, and the $i$-th term of each row is the number of $i$'s in the next row, listed in increasing order:
\[
  (1,1),\quad (1,2),\quad (1,2,2),\quad (1,2,2,3,3),\quad (1,2,2,3,3,4,4,4,5,5,5),\ \dots
\]
The final terms of the rows form Levine's sequence A014644 \cite{A014644},
\[
  A = 1,2,2,3,5,11,38,272,6474,1090483,4363282578,\dots
\]
(terms $A(1),\dots,A(11)$ were computed by C.~Mallows and $A(12),\dots,A(15)$ by J.~Claes in 2005; the b-file of the entry now gives nineteen terms, the last four computed by Claes and Miyamoto \cite{A014644,CM}).

The OEIS entry A113676 \cite{A113676}, by Floor van Lamoen and Paul D.\ Hanna (November~6, 2005), takes Golomb's sequence with one $2$ removed, $1,2,3,3,4,4,4,5,5,5,\dots$, and cuts it into consecutive rows: the first row is $(1)$, and the number of terms of row $n$ is the sum of row $n-1$. The triangle begins
\[
  (1);\quad (2);\quad (3,3);\quad (4,4,4,5,5,5);\ \dots
\]
and $a(n)$ is the number of terms of row $n$: $1,1,2,6,27,234,6202,1084009,4362192095$.
The entry records two conjectures, quoted in Section~\ref{sec:statement}: the row lengths $a(n)$, $n\ge2$, are the first differences of $A$ from its third term on, and the final terms of the rows are the terms of $A$ with its duplicate $2$ removed. We prove both (Theorem~\ref{thm:main}).

\subsection*{Method}
Write $P(l)=(G(1),\dots,G(l))$ and $L_k$ for the number of terms of row $k$ of Levine's array.
Since $G$ describes itself, the Levine step sends $P(l)$ to the list $1^{G(1)}2^{G(2)}\cdots l^{G(l)}$, which is again an initial segment of $G$; so row $k\ge2$ of Levine's array is $P(L_k)$ (Lemma~\ref{lem:levine}). Taking sums gives the telescoping identity $G(L_k+1)+\dots+G(L_{k+1})=L_{k+2}-L_{k+1}$ (Lemma~\ref{lem:tele}), and an induction shows that row $n\ge2$ of the triangle is exactly the block $G(L_n+1),\dots,G(L_{n+1})$ (Lemma~\ref{lem:rowend}). Since $L_k=A(k+1)$ and the last term of $P(L_{n+1})$ is $A(n+1)$, both conjectures follow.

\subsection*{Status and prior work}
The entry A113676 (revision \#9, March~5, 2025) still lists both statements as conjectures, with no proof or comment on them.
The comments of A001462 contain a related conjecture of van Lamoen of the same date, phrased in terms of a differently described triangle and partial sums of its row sums; it is still worded as a conjecture in the current revision, and we do not address it here.
The key ingredient of our proof is known: the note ``Rows converge to A001462'' in A014643 reflects it, and Claes and Miyamoto \cite[p.~3, eqs.~(3)--(4)]{CM} observe that the Levine step (their ``golombic operator'' $\Gamma$) maps the first $n$ terms of Golomb's sequence to its first A001463$(n)$ terms \cite{A001463}, so that the rows of A014643 are initial segments of $G$, and they record how the length, the sum and the last term of a list relate to those of its image under the Levine step (our Lemma~\ref{lem:gamma}). Their paper does not mention A113676 or the triangle, and does not derive the two statements proved here.
The math0 referees also checked A014644, A014643, A011784, A012257, A001463, A113722, A113724, A113730 and A304679, and web searches (including arXiv and the SeqFan list) found no resolution.
We therefore classify the result as new, with the caveat that it is a short consequence (a telescoping sum and an induction) of the initial-segment property above.

\section{Statement}\label{sec:statement}

For a finite list $c=(c_1,\dots,c_l)$ of natural numbers we write $\len c=l$ for its length, $\Sigma c=c_1+\dots+c_l$ for its sum and $\last(c)=c_l$ for its last term (if $l\ge1$). We write $i^{m}$ for the list of $m$ copies of $i$ and juxtaposition for concatenation, and define the \emph{Levine step}
\[
  \Gamma(c_1,\dots,c_l) = 1^{c_1}\,2^{c_2}\cdots l^{c_l}.
\]

\begin{definition}\label{def:main}
Let $G$ be Golomb's sequence and $P(l)=(G(1),\dots,G(l))$.
\begin{enumerate}
\item Let $H(1)=1$ and $H(j)=G(j+1)$ for $j\ge2$, so $H=1,2,3,3,4,4,4,5,5,5,\dots$.
\item Let $E(0)=0$, $E(1)=1$ and $E(n+2)=E(n+1)+\sum_{E(n)<j\le E(n+1)}H(j)$. Row $n\ge1$ of the triangle is $R_n=(H(j))_{E(n-1)<j\le E(n)}$, and $a(n)=\len{R_n}$. Thus $R_1=(1)$ and $\len{R_{n+2}}=\Sigma R_{n+1}$.
\item Let $\Lev_1=(1,1)$ and $\Lev_{k+1}=\Gamma(\Lev_k)$ for $k\ge1$; let $A(k)=\last(\Lev_k)$ and $L_k=\len{\Lev_k}$.
\end{enumerate}
\end{definition}

The sequence $E$ is nondecreasing, so $a(n)=E(n)-E(n-1)$. One has $\Lev_2=(1,2)$, $\Lev_3=(1,2,2)$, so $L_1=L_2=2$ and $L_3=3$.

The COMMENTS section of A113676 \cite{A113676} (entry by Floor van Lamoen and Paul D.\ Hanna, Nov~06 2005) contains, verbatim:
\begin{quote}
{\small\texttt{Conjecture: a(n) for n>1 gives first differences of Lionel Levine's sequence A014644(n) for n>=3.}}\\
{\small\texttt{Conjecture: Final elements of the rows form A014644 except for duplicate 2.}}
\end{quote}
The name of the entry refers to ``A001642''; this is a typo for A001462, as the example rows of the entry show.

\begin{theorem}\label{thm:main}
\textup{(1)} For every $n\ge2$, $a(n)=A(n+2)-A(n+1)$.
\textup{(2)} $\last(R_1)=A(1)$, and $\last(R_n)=A(n+1)$ for every $n\ge2$.
\end{theorem}

The proof in Section~\ref{sec:proof} gives slightly more, which is not part of the formal statement below (it appears only as intermediate steps of the Lean proof): for $n\ge2$ row $n$ of the triangle is the block $R_n=(G(i))_{A(n+1)<i\le A(n+2)}$ of Golomb's sequence, and $G(A(n+2))=A(n+1)$.

The indexing in (1) is the one forced by the data: $a(2)=1=A(4)-A(3)$ and $a(3)=2=A(5)-A(4)$, whereas the shift $A(n+1)-A(n)$ would give $a(3)=1$. In (2), the final terms $1,2,3,5,11,38,272,\dots$ of the rows are $A=1,2,2,3,5,11,38,272,\dots$ with one of the two equal terms $A(2)=A(3)=2$ removed.

\subsection*{Lean formalization}
The statement was formalized by math0 in \path{lean/Math0/Statements/OEIS/A113676.lean}, reviewed for fidelity and frozen (SHA-256 \texttt{a80c3eb1\ldots afab5}) before any proof attempt. Its definitions and the statement are, verbatim (the import, docstrings, section headers, blank lines and all theorems omitted; \texttt{List.range' s c} is $[s,s+1,\dots,s+c-1]$):
\begin{lstlisting}
namespace Math0.Statements.A113676
def golombAux : ℕ → List ℕ
  | 0 => [1, 2, 2]
  | m + 1 => golombAux m ++ List.replicate ((golombAux m).getD (m + 2) 0) (m + 3)
def golomb : ℕ → ℕ
  | 0 => 0
  | n + 1 => (golombAux n).getD n 0
def H (j : ℕ) : ℕ := if j ≤ 1 then golomb j else golomb (j + 1)
def rowEnd : ℕ → ℕ
  | 0 => 0
  | 1 => 1
  | n + 2 => rowEnd (n + 1) +
      ((List.range' (rowEnd n + 1) (rowEnd (n + 1) - rowEnd n)).map H).sum
def row (n : ℕ) : List ℕ :=
  (List.range' (rowEnd (n - 1) + 1) (rowEnd n - rowEnd (n - 1))).map H
def a (n : ℕ) : ℕ := (row n).length
def levineAux : List ℕ → ℕ → List ℕ
  | [], _ => []
  | c :: l, i => List.replicate c i ++ levineAux l (i + 1)
def levineRow : ℕ → List ℕ
  | 0 => []
  | 1 => [1, 1]
  | k + 2 => levineAux (levineRow (k + 1)) 1
def A014644 (k : ℕ) : ℕ := (levineRow k).getLastD 0
def statement : Prop :=
  (∀ n : ℕ, 2 ≤ n → (a n : ℤ) = (A014644 (n + 2) : ℤ) - (A014644 (n + 1) : ℤ)) ∧
  ((row 1).getLast? = some (A014644 1) ∧
    ∀ n : ℕ, 2 ≤ n → (row n).getLast? = some (A014644 (n + 1)))
end Math0.Statements.A113676
\end{lstlisting}
The encoding matches Definition~\ref{def:main}: \decl{rowEnd} is $E$, \decl{row} is $R_n$, \texttt{levineAux l 1} is $\Gamma(l)$, \decl{levineRow} is $\Lev_k$ and \decl{A014644} is $A$ (all with offset $1$; the values at $0$ are never used). Golomb's sequence is computed by its self-generating construction: \texttt{golombAux m} starts from $1,2,2$ and appends, for $k=3,\dots,m+2$, the number $k$ repeated $G(k)$ times, where $G(k)$ is read off the list built so far. The file proves (proofs omitted)
\begin{lstlisting}
theorem golombAux_eq_flatMap (m : ℕ) :
    golombAux m = (List.range' 1 (m + 2)).flatMap (fun k => List.replicate (golomb k) k)
theorem golombAux_eq_map (m : ℕ) :
    golombAux m = (List.range' 1 (golombAux m).length).map golomb
theorem a_data : (List.range 9).map (fun i => a (i + 1)) =
    [1, 1, 2, 6, 27, 234, 6202, 1084009, 4362192095]
theorem A014644_data : (List.range 11).map (fun i => A014644 (i + 1)) =
    [1, 2, 2, 3, 5, 11, 38, 272, 6474, 1090483, 4363282578]
\end{lstlisting}
The first two say that every stage is $P(\ell)=1^{G(1)}2^{G(2)}\cdots(m+2)^{G(m+2)}$ for its length $\ell$, which is the self-describing property of A001462; together with $G(1)=1$ it determines the sequence. The file also checks by kernel evaluation all $84$ terms of A001462 listed in the OEIS (\decl{golomb_data}), the four example rows of A113676 (\decl{rows_example}), and the last terms $1,2,3,5,11,38,272$ of rows $1$ to $7$ (\decl{rowLast_data}); \decl{a_data} and \decl{A014644_data} check all $9$ listed terms of A113676 and the first $11$ of A014644.
The large values are certified without computing long lists in the kernel. The iterates of $\Gamma$ starting from $(1,2,2)$ are proved to be initial segments of $G$ (\decl{golombFast_eq_golombAux}); the fifth one, $P(6474)$, is evaluated by \texttt{decide +kernel}, and the sixth one, $P(1090483)$, is never built: its length and sum are obtained from the sum $1090483$ and the weighted sum $\sum_{i\le 6474}iG(i)=4363282578$ of $P(6474)$, via the lemmas that $\len{\Gamma(c)}=\Sigma c$ and $\Sigma\Gamma(c)=\sum_i i\,c_i$. The rows $1,\dots,7$ and $a(1),\dots,a(8)$ are read off $P(6474)$; $a(9)=\sum_{6475\le i\le 1090483}G(i)=4363282578-1090483$ is obtained as above; and the same two lemmas applied to $\Lev_8=P(6474)$, together with $A(k+2)=L_{k+1}$ (\decl{A014644_add_two}), give $A(9),A(10),A(11)$. These checks only concern the fidelity of the definitions; the proof of Theorem~\ref{thm:main} uses none of them.

The statement is computed in $\Z$, so the difference in (1) is a true difference, and ``the last term of row $n$ is $x$'' is \texttt{(row n).getLast? = some x}, which fails for an empty row.

\section{Proof}\label{sec:proof}

The proof follows the Lean file \path{lean/Math0/Frontier/OEIS/A113676.lean}; the names of the corresponding Lean lemmas are given in brackets.

\begin{lemma}[Levine step]\label{lem:gamma}
Let $c=(c_1,\dots,c_l)$. Then $\len{\Gamma(c)}=\Sigma c$ {\upshape[\decl{length_levineAux}]}, and all terms of $\Gamma(c)$ are $\ge1$ {\upshape[\decl{levineAux_pos}]}. If $l\ge1$ and all $c_i\ge1$, then $\last(\Gamma(c))=l$ {\upshape[\decl{getLast?_levineAux}]}. If all $c_i\ge1$, then $\Sigma c\ge l$.
\end{lemma}

\begin{proof}
The first two claims are immediate from the definition. If $l\ge1$ and $c_l\ge1$, the last block $l^{c_l}$ of $\Gamma(c)$ is nonempty, so $\Gamma(c)$ ends in $l$. The last claim is clear.
\end{proof}

\begin{lemma}[Golomb's sequence]\label{lem:golomb}
$G(i)\ge1$ for $i\ge1$ {\upshape[\decl{golomb_pos}]}, and for every $l\ge2$,
\[
  \Gamma(P(l))=1^{G(1)}2^{G(2)}\cdots l^{G(l)}=P(l'),\qquad l'=G(1)+\dots+G(l).
\]
\end{lemma}

\begin{proof}
Positivity holds because every stage of the construction consists of positive numbers [\decl{golombAux_pos_length}]. Let $m=l-2$. By \decl{golombAux_eq_flatMap}, the stage \texttt{golombAux m} is $1^{G(1)}\cdots l^{G(l)}=\Gamma(P(l))$, and by \decl{golombAux_eq_map} it is $P(l')$ for its length $l'$, which is $\Sigma P(l)$ by Lemma~\ref{lem:gamma}.
\end{proof}

\begin{lemma}[Levine's rows]\label{lem:levine}
For $k\ge1$: $L_k\ge2$ {\upshape[\decl{two_le_length_levineRow}]}, $L_{k+1}=\Sigma\Lev_k\ge L_k$ {\upshape[\decl{length_levineRow_succ}, \decl{length_levineRow_le}]}, and $A(k+1)=L_k$ {\upshape[\decl{A014644_add_two}]}. For $k\ge2$: $\Lev_k=P(L_k)$ {\upshape[\decl{levineRow_eq}]}.
\end{lemma}

\begin{proof}
All terms of $\Lev_k$ are $\ge1$ ($\Lev_1=(1,1)$ and Lemma~\ref{lem:gamma}). Hence $L_{k+1}=\len{\Gamma(\Lev_k)}=\Sigma\Lev_k\ge L_k$, and $L_k\ge L_1=2$ by induction. Since $\Lev_k$ is nonempty with positive terms, $A(k+1)=\last(\Gamma(\Lev_k))=L_k$ by Lemma~\ref{lem:gamma}.
For the last claim we induct on $k\ge2$. For $k=2$, $\Lev_2=(1,2)=P(2)$ (checked by \texttt{decide}). If $\Lev_k=P(L_k)$ with $L_k\ge2$, then by Lemma~\ref{lem:golomb} $\Lev_{k+1}=\Gamma(P(L_k))=P(l')$ for some $l'$, and necessarily $l'=L_{k+1}$.
\end{proof}

\begin{lemma}[Telescoping]\label{lem:tele}
For $k\ge2$ {\upshape[\decl{sum_block}]},
\[
  L_{k+1}+\sum_{i=L_k+1}^{L_{k+1}}G(i)=L_{k+2}.
\]
\end{lemma}

\begin{proof}
As $L_k\le L_{k+1}$, $P(L_{k+1})$ is $P(L_k)$ followed by $(G(L_k+1),\dots,G(L_{k+1}))$. By Lemma~\ref{lem:levine}, $P(L_{k+1})=\Lev_{k+1}$ and $P(L_k)=\Lev_k$; taking sums and using $\Sigma\Lev_{k+1}=L_{k+2}$, $\Sigma\Lev_k=L_{k+1}$ gives the claim.
\end{proof}

\begin{lemma}[Row ends]\label{lem:rowend}
For every $n\ge1$, $E(n)+1=L_{n+1}$ and $E(n+1)+1=L_{n+2}$ {\upshape[\decl{rowEnd_eq}]}.
\end{lemma}

\begin{proof}
Induction on $n$. For $n=1$: $E(1)=1$, $E(2)=E(1)+H(1)=2$, $L_2=2$ and $L_3=3$ (checked by \texttt{decide}). Assume both equalities for $n$; the first one for $n+1$ is the second one for $n$. By definition,
\[
  E(n+2)=E(n+1)+\sum_{j=E(n)+1}^{E(n+1)}H(j).
\]
Since $E(n)+1=L_{n+1}\ge2$, every index $j$ in the sum is $\ge2$, so $H(j)=G(j+1)$ [\decl{map_H_range'_golomb}], and the sum equals $\sum_{i=L_{n+1}+1}^{L_{n+2}}G(i)$, which is $L_{n+3}-L_{n+2}$ by Lemma~\ref{lem:tele} with $k=n+1\ge2$. Hence $E(n+2)+1=L_{n+2}+(L_{n+3}-L_{n+2})=L_{n+3}$.
\end{proof}

\begin{lemma}[Strict growth]\label{lem:strict}
$L_{k+1}<L_{k+2}$ for every $k\ge1$ {\upshape[\decl{length_levineRow_lt}]}.
\end{lemma}

\begin{proof}
For $k=1$, $L_2=2<3=L_3$. If $L_{k+1}<L_{k+2}$, then by Lemma~\ref{lem:tele} with $k+1$, $L_{k+3}-L_{k+2}$ is a sum of $L_{k+2}-L_{k+1}\ge1$ Golomb terms, each $\ge1$ by Lemma~\ref{lem:golomb} [\decl{le_sum_map_golomb}], so $L_{k+3}-L_{k+2}\ge1$.
\end{proof}

\begin{proof}[Proof of Theorem~\ref{thm:main}]
Row $1$ is $(1)$ [\decl{row_one}] and $\Lev_1=(1,1)$, so $\last(R_1)=1=A(1)$.
Let $n\ge2$. By Lemma~\ref{lem:rowend} applied to $n-1\ge1$, $E(n-1)+1=L_n\ge2$ and $E(n)+1=L_{n+1}$; as in the proof of Lemma~\ref{lem:rowend}, $H(j)=G(j+1)$ for $j>E(n-1)$, so
\[
  R_n=(H(j))_{E(n-1)<j\le E(n)}=(G(i))_{L_n<i\le L_{n+1}}.
\]
By Lemma~\ref{lem:levine}, $L_n=A(n+1)$ and $L_{n+1}=A(n+2)$, which gives the description of $R_n$ stated after the theorem.

(1) $a(n)=\len{R_n}=L_{n+1}-L_n=A(n+2)-A(n+1)$; this is a true difference, since $L_n\le L_{n+1}$.

(2) By Lemma~\ref{lem:strict}, $L_n<L_{n+1}$, so $R_n$ is nonempty and $\last(R_n)=G(L_{n+1})$. By Lemma~\ref{lem:levine}, $\Lev_{n+1}=P(L_{n+1})$ with $L_{n+1}\ge1$, so $A(n+1)=\last(\Lev_{n+1})=G(L_{n+1})$. Hence $\last(R_n)=A(n+1)$, and $G(A(n+2))=A(n+1)$.
\end{proof}

\begin{remark}
Lemmas~\ref{lem:golomb} and~\ref{lem:levine} also give $A(k+2)=L_{k+1}=\Sigma P(L_k)=\mathrm{A001463}(A(k+1))$ for $k\ge2$, where A001463 is the sequence of partial sums of $G$ \cite{A001463}; this identity, which is not stated separately in the Lean file, is a restatement of the observation of Claes and Miyamoto \cite[p.~3]{CM} that $\Gamma$ maps $P(n)$ to $P(\mathrm{A001463}(n))$. For example, $L_2,L_3,\dots=2,3,5,11,38,272,\dots$, and the rows $R_2=(G(3))=(2)$, $R_3=(G(4),G(5))=(3,3)$ and $R_4=(G(6),\dots,G(11))=(4,4,4,5,5,5)$ are the blocks of $G$ between consecutive values of $L_k$.
\end{remark}

\section{Machine verification}\label{sec:lean}

\subsection*{Files and declarations}
The proof is in the file \path{lean/Math0/Frontier/OEIS/A113676.lean} of the math0 repository, \url{https://github.com/Foundation-0/math0}. It imports the frozen statement file \path{lean/Math0/Statements/OEIS/A113676.lean} and the audit commands in \path{lean/Math0/Audit.lean}. The main declaration is \decl{Math0.Frontier.A113676.result}, of type \decl{Math0.Statements.A113676.statement}. The private lemmas \decl{golomb_pos}, \decl{le_sum_map_golomb}, \decl{two_le_length_levineRow}, \decl{length_levineRow_succ}, \decl{length_levineRow_le}, \decl{levineRow_eq}, \decl{sum_block}, \decl{rowEnd_eq}, \decl{length_levineRow_lt} and \decl{getLast?_map_range'} correspond to Lemmas~\ref{lem:golomb}--\ref{lem:strict} and to the last step of the proof of Theorem~\ref{thm:main}. From the statement file the proof uses general lemmas such as \decl{golombAux_eq_flatMap}, \decl{golombAux_eq_map}, \decl{levineAux_map}, \decl{length_levineAux}, \decl{A014644_add_two} and \decl{map_H_range'_golomb}, but none of the data checks. The proof file contains no \texttt{sorry}, no \texttt{native\_decide} and no \texttt{set\_option}; the tactic \texttt{decide} is used only for the base cases about $\Lev_1,\Lev_2,\Lev_3$, $P(2)$, $E(1)$ and $E(2)$, which the kernel checks by evaluation without additional axioms.

\subsection*{Audit}
The file ends with the command
\begin{lstlisting}
#math0_settles Math0.Frontier.A113676.result Math0.Statements.A113676.statement
\end{lstlisting}
Defined in \path{lean/Math0/Audit.lean} for statements that math0 formalizes itself, it fails unless (1) \decl{statement} is a \texttt{Prop}-valued definition without universe parameters; (2) \decl{statement} and \decl{result} use no axioms besides \decl{propext}, \decl{Classical.choice} and \decl{Quot.sound} (rejecting \decl{sorryAx}, \decl{Lean.ofReduceBool} from \texttt{native\_decide}, and custom axioms); and (3) the type of \decl{result} is syntactically the statement or its negation, up to metadata, with no unfolding. Here it reports
\begin{lstlisting}[basicstyle=\ttfamily\scriptsize]
math0 audit: 'Math0.Statements.A113676.statement' depends on axioms [propext]
math0 audit: 'Math0.Frontier.A113676.result' depends on axioms [Classical.choice, Quot.sound, propext]
math0 statement check: 'Math0.Frontier.A113676.result' proves 'Math0.Statements.A113676.statement'
\end{lstlisting}

\subsection*{Toolchain and reproduction}
We use Lean~4 \cite{Lean4} \texttt{v4.33.1} and Mathlib \cite{Mathlib} \texttt{v4.33.1}, pinned in \path{lean/}; the repository also pins formal-conjectures \cite{FC} at commit \texttt{c77fdd4}, which this result does not use, since the statement is not part of that repository.
To reproduce, run \path{scripts/verify.sh} from the repository root: it fetches the Mathlib build cache, builds all math0 results with \texttt{lake build} (re-running every audit command) and checks that the pinned \texttt{formal\_conjectures} and \texttt{mathlib} checkouts are unmodified. The single file can be checked with \texttt{lake env lean Math0/Frontier/OEIS/A113676.lean} in the directory \path{lean}, after \texttt{lake build Math0.Statements.OEIS.A113676}.

\begin{thebibliography}{99}

\bibitem{CM}
J.~Claes and R.~Miyamoto, \emph{Golombic and Levine sequences}, arXiv:2602.10992 [math.CO], 2026.

\bibitem{FC}
M.~Firsching, P.~Lezeau, S.~Mercuri, M.~Z. Horv\'ath, Y.~Dillies, C.~S\"onne, E.~Wieser, F.~Zhang, T.~Hubert, B.~Ag\"uera y Arcas, and P.~Kohli, \emph{Formal Conjectures: An open and evolving benchmark for verified discovery in mathematics}, arXiv:2605.13171, 2026. Repository: \url{https://github.com/google-deepmind/formal-conjectures}.

\bibitem{Go}
S.~W. Golomb, \emph{Problem 5407}, Amer. Math. Monthly \textbf{73} (1966), no.~6, 674.

\bibitem{Guy}
R.~K. Guy, \emph{Unsolved Problems in Number Theory}, 3rd ed., Springer, 2004, Section E25, 347--348.

\bibitem{Mathlib}
The mathlib Community, \emph{The Lean mathematical library}, in: Proceedings of the 9th ACM SIGPLAN International Conference on Certified Programs and Proofs (CPP 2020), ACM, 2020, 367--381.

\bibitem{Lean4}
L.~de Moura and S.~Ullrich, \emph{The Lean 4 theorem prover and programming language}, in: Automated Deduction -- CADE 28, Lecture Notes in Comput. Sci. \textbf{12699}, Springer, 2021, 625--635.

\bibitem{A001462}
The OEIS Foundation Inc., \url{https://oeis.org/A001462} (Golomb's sequence), entry by N.~J.~A. Sloane.

\bibitem{A001463}
The OEIS Foundation Inc., \url{https://oeis.org/A001463} (partial sums of A001462), entry by N.~J.~A. Sloane.

\bibitem{A014643}
The OEIS Foundation Inc., \url{https://oeis.org/A014643} (Levine's array), entry by N.~J.~A. Sloane.

\bibitem{A014644}
The OEIS Foundation Inc., \url{https://oeis.org/A014644} (Levine's sequence), entry by N.~J.~A. Sloane.

\bibitem{A113676}
The OEIS Foundation Inc., \url{https://oeis.org/A113676}, entry and conjectures by F.~van Lamoen and P.~D. Hanna (Nov.~6, 2005); revision \#9, Mar.~5, 2025.

\end{thebibliography}
\end{document}
